\section*{Bab \ref{ch:b2:diatomic-mos} --- Orbital Molekul Molekul Diatomik}

\begin{solution}{exo:b2:diatomic-mos:1}
$E_+ = (-13.6 - 9.0)/1.60 = \qty{-14.1}{eV}$; $E_- = (-13.6 + 9.0)/0.40 = \qty{-11.5}{eV}$.
Dua elektron dalam $\psi_+$: $-28.25$ dibandingkan
$2 \times (-13.6) = \qty{-27.2}{eV}$, terikat sebesar \qty{1.05}{eV} dalam
model kasar ini.
\end{solution}

\begin{solution}{exo:b2:diatomic-mos:2}
\ce{He2}: $(2 - 2)/2 = 0$; \ce{He2+}: $(2 - 1)/2 = 0.5$; \ce{Li2}: 1; \ce{B2}:
enam elektron valensi, $\sigma_{2s}^2\sigma_{2s}^{*2}\pi_{2p}^2$, \omterm{def:b2:diatomic-mos:bond-order}{orde ikatan} 1.
\end{solution}

\begin{solution}{exo:b2:diatomic-mos:3}
\ce{B2} (dua elektron tak berpasangan dalam kedua orbital $\pi_{2p}$, dengan
pencampuran s--p) dan \ce{O2} (dua dalam $\pi^*$). \ce{C2}, \ce{N2}, \ce{F2}
\omterm{def:b2:diatomic-mos:magnetism}{diamagnetik}.
\end{solution}

\begin{solution}{exo:b2:diatomic-mos:4}
Tak nol: (1s, $2\mathrm p_z$), ($2\mathrm p_x$, $2\mathrm p_x$), (2s,
$2\mathrm p_z$). Nol karena simetri: (1s, $2\mathrm p_x$) dan ($2\mathrm p_x$,
$2\mathrm p_y$).
\end{solution}

\begin{solution}{exo:b2:diatomic-mos:5}
Ionisasi mengambil elektron ikatan: \omterm{def:b2:diatomic-mos:bond-order}{orde ikatan} 2.5 alih-alih 3, sehingga
ikatannya lebih panjang dan lebih lemah, seperti yang ditunjukkan oleh
$D_0(\ce{N2+}) = \qty{8.71}{eV} < D_0(\ce{N2}) = \qty{9.76}{eV}$.
\end{solution}

\begin{solution}{exo:b2:diatomic-mos:6}
Orbital-orbital oksigen terletak lebih rendah (oksigen lebih
elektronegatif): menurut \cref{thm:b2:diatomic-mos:heteronuclear}, \omterm{def:b2:diatomic-mos:bonding}{orbital
ikatan} berbobot pada atom yang lebih rendah, oksigen, dan orbital-orbital
yang lebih tinggi berbobot pada karbon. Orbital terisi tertinggi, sebuah
orbital $\sigma$ yang sebagian besar bersifat karbon dan mengarah menjauhi
oksigen, adalah pasangan elektron bebas pada karbon.
\end{solution}

\begin{solution}{exo:b2:diatomic-mos:7}
Hanya $2\mathrm p_z$ fluor (sepanjang sumbu) yang mempunyai simetri 1s
hidrogen: keduanya membentuk orbital $\sigma$ yang berbobot pada F dan
orbital $\sigma^*$ yang berbobot pada H. Orbital $2\mathrm p_x$, $2\mathrm p_y$
fluor (tumpang-tindih nol) dan 2s-nya yang rendah (terlalu jauh energinya)
tetap nonikatan: tiga pasangan elektron bebas. Delapan elektron valensi
mengisi 2s, $\sigma$, dan kedua orbital p nonikatan.
\end{solution}

\begin{solution}{exo:b2:diatomic-mos:8}
$\bar\alpha = \qty{-12}{eV}$, $\Delta = \qty{4}{eV}$, $\sqrt{4 + 9} = \qty{3.61}{eV}$:
$E = \qty{-15.6}{eV}$ dan \qty{-8.4}{eV}. Untuk tingkat yang lebih rendah
$(\alpha_B - E)c_B = -\beta
c_A$: $1.61c_B = 3.0c_A$, $c_A/c_B = 0.535$,
$c_B^2 = 1/(1 + 0.286) = 0.78$.
\end{solution}

\begin{solution}{exo:b2:diatomic-mos:9}
Delapan elektron valensi: $\sigma_{2s}^2\sigma_{2s}^{*2}\pi_{2p}^4$, dengan
orbital $\sigma_{2p}$ (yang terdorong ke atas $\pi_{2p}$ oleh pencampuran)
kosong. \omterm{def:b2:diatomic-mos:bond-order}{Orde ikatan} $(6 - 2)/2 = 2$, dari kedua orbital $\pi$ yang terisi.
\end{solution}

\begin{solution}{exo:b2:diatomic-mos:10}
NO, sebelas elektron valensi, satu dalam $\pi^*$: \omterm{def:b2:diatomic-mos:bond-order}{orde ikatan} 2.5.
\ce{NO+}, sepuluh: \omterm{def:b2:diatomic-mos:bond-order}{orde ikatan} 3.
$D_0(\ce{NO+}) = 6.51 + 13.62 - 9.26 = \qty{10.87}{eV}$: jauh lebih kuat
daripada NO, karena elektron yang diambil bersifat antiikatan.
\end{solution}

\begin{solution}{exo:b2:diatomic-mos:11}
$E_+ + E_- = \dfrac{(\alpha + \beta)(1 - S) + (\alpha - \beta)(1 + S)}{1 - S^2} = \dfrac{2(\alpha -
\beta S)}{1 - S^2}$. Nilai ini melampaui $2\alpha$ jika
$\alpha - \beta S > \alpha - \alpha S^2$, yaitu $\beta < \alpha S$ (dengan
membagi oleh $-S < 0$). Dengan empat elektron, dua pada setiap orbital,
energi totalnya di atas energi orbital-orbital yang terpisah: dua orbital
yang terisi saling menolak.
\end{solution}

\begin{solution}{exo:b2:diatomic-mos:12}
\omterm{def:b2:diatomic-mos:bond-order}{Orde ikatan} 2.5, 2, 1.5, 1: menurut urutan ini ikatannya makin panjang dan
makin lemah. Hidrogen peroksida mengandung ikatan tunggal O--O peroksida
(satuan \ce{O2^2-}); \ce{KO2} mengandung ion superoksida \ce{O2-}, yang
\omterm{def:b2:diatomic-mos:magnetism}{paramagnetik}.
\end{solution}

\begin{solution}{pb:b2:diatomic-mos:1}
\textbf{1.} Dua belas.
\textbf{2.} $\sigma_{2s}^2\,\sigma_{2s}^{*2}\,\sigma_{2p}^2\,\pi_{2p}^4\,\pi_{2p}^{*2}$,
dengan kedua elektron $\pi^*$ pada orbital yang berbeda.
\textbf{3.} $(8 - 4)/2 = 2$.
\textbf{4.} Dua: \ce{O2} \omterm{def:b2:diatomic-mos:magnetism}{paramagnetik}.
\textbf{5.} Struktur itu memasangkan semua elektron, padahal dua di
antaranya menempati sendiri-sendiri dua orbital $\pi^*$ yang
terdegenerasi.
\textbf{6.} $D_0 = 2 \times 246.79 = \qty{493.6}{kJ/mol}$, \qty{5.12}{eV}.
\textbf{7.} \ce{O2+}: $\pi^{*1}$, \omterm{def:b2:diatomic-mos:bond-order}{orde ikatan} 2.5.
\textbf{8.} \ce{O2-}: $\pi^{*3}$, \omterm{def:b2:diatomic-mos:bond-order}{orde ikatan} 1.5.
\textbf{9.} \ce{O2^2-}: $\pi^{*4}$, \omterm{def:b2:diatomic-mos:bond-order}{orde ikatan} 1.
\textbf{10.} \ce{O2+} satu, \ce{O2-} satu, \ce{O2^2-} tidak ada.
\textbf{11.} \ce{O2+} $<$ \ce{O2} $<$ \ce{O2-} $<$ \ce{O2^2-}.
\textbf{12.} Ion peroksida.
\textbf{13.} \ce{O2+}, \ce{O2}, dan \ce{O2-}.
\textbf{14.} Dari sebuah orbital $\pi^*$.
\textbf{15.} Orbital $\pi^*$, yang bersifat antiikatan, terletak di atas
tingkat 2p atom: elektronnya kurang terikat.
\textbf{16.} Dua jalur dari \ce{O2} ke $\ce{O} + \ce{O+} + \mathrm e^-$:
mendisosiasi lalu mengionisasi sebuah atom, $D_0(\ce{O2}) + I(\ce{O})$; atau
mengionisasi molekul lalu mendisosiasi ionnya, $I(\ce{O2}) + D_0(\ce{O2+})$.
Keduanya sama: $D_0(\ce{O2+}) = 5.12 + 13.62 - 12.07 = \qty{6.66}{eV}$.
\textbf{17.} Lebih besar daripada $D_0(\ce{O2}) = \qty{5.12}{eV}$: mengambil
elektron antiikatan memperkuat ikatan (orde 2.5).
\textbf{18.} $D_0(\ce{N2}) = 2 \times 470.82/96.485 = \qty{9.76}{eV}$; $D_0(\ce{N2+}) = 9.76 +
14.53 - 15.58 = \qty{8.71}{eV}$.
\textbf{19.} Orbital terisi tertinggi \ce{N2} bersifat ikatan, di bawah
tingkat 2p atom.
\textbf{20.} Ketika elektron yang diambil berasal dari \omterm{def:b2:diatomic-mos:bonding}{orbital antiikatan}.
\textbf{21.} Aturan Hund (jumlah maksimum spin sejajar dalam orbital-orbital
yang terdegenerasi).
\textbf{22.} Kedua elektron $\pi^*$ pada orbital yang sama, dengan spin
berpasangan.
\textbf{23.} Dua elektron dalam orbital yang sama lebih kuat saling menolak,
dan kehilangan stabilisasi pertukaran dari spin yang sejajar.
\textbf{24.} \omterm{def:b2:diatomic-mos:bond-order}{Orde ikatan} tetap 2; semua elektron berpasangan: \omterm{def:b2:diatomic-mos:magnetism}{diamagnetik}.
\textbf{25.} \ce{O2+} mempunyai \omterm{def:b2:diatomic-mos:bond-order}{orde ikatan} $\boldsymbol{2.5}$ dan energi
disosiasi $\boldsymbol{\approx \qty{6.66}{eV}}$, dibandingkan dengan 2 dan
\qty{5.12}{eV} untuk \ce{O2}.
\end{solution}
